\section{Hacia la inteligencia experiencial: de Context Engineering a Context Learning}

Ma Xiaoteng (马骁腾) llama a la siguiente etapa de la capacidad de los Agent \textbf{Experiential Intelligence}: el sistema no solo ensambla el context en una única solicitud, sino que aprende, a partir de la interacción de largo plazo, «qué información resulta más útil para este usuario, esta tarea y este momento». Esto exige elevar el context engineering de reglas manuales a un context learning entrenable, evaluable y sostenible en servicio.

\sourcefigure{0.93}{lecture06/slides-images/slide-002-00030s.jpg}{Inteligencia experiencial: que el sistema obtenga capacidades personalizadas a partir de la interacción de largo plazo}{00:00:30--00:01:00}

\subsection{Por qué no basta con ampliar solo la escala del modelo}

La escala del modelo sigue creciendo, pero lo que cada usuario realmente necesita son preferencias personales estables, el estado del proyecto y la experiencia histórica. Los parámetros generales no pueden absorber a tiempo información privada y en constante cambio; y meter todo el historial en el prompt trae consigo costo de contexto, ruido y problemas de privacidad.

\sourcefigure{0.93}{lecture06/slides-images/slide-003-00060s.jpg}{La escala del modelo y la escala de la experiencia individual no siguen la misma curva de escalabilidad}{00:01:00--00:01:45}

\sourcefigure{0.93}{lecture06/slides-images/slide-004-00105s.jpg}{La «segunda mitad» de la IA pasa de la capacidad general a la experiencia y la personalización}{00:01:45--00:02:00}

\begin{importantbox}{El objeto del Context Learning}
Lo que se aprende no es un prompt fijo, sino la estrategia de «cómo seleccionar, comprimir, ordenar y actualizar el context a partir de la memoria, el entorno, el estado del usuario y los resultados de las herramientas». Esta estrategia decide qué experiencias entran al razonamiento actual, cuáles se escriben de vuelta al estado de largo plazo y cuáles deben olvidarse.
\end{importantbox}

\subsection{LoRA as Memory: un portador de experiencia eficiente en parámetros}

El informe explora tratar el LoRA adapter como memoria entrenable. A diferencia del fine-tuning completo, LoRA solo aprende un incremento de bajo rango

$$
W'=W+BA,\qquad \mathrm{rank}(BA)=r\ll d,
$$

donde el peso base $W$ permanece congelado y $A,B$ almacenan la experiencia específica. Distintos usuarios o tareas pueden tener adapters independientes, que se cargan según se necesiten al momento de servir.

\sourcefigure{0.92}{lecture06/slides-images/slide-014-00510s.jpg}{El límite entre context engineering y context learning}{00:08:30--00:11:30}

\sourcefigure{0.92}{lecture06/slides-images/slide-015-00690s.jpg}{LoRA as Memory: portar experiencia de largo plazo mediante incrementos de parámetros}{00:11:30--00:13:15}

\begin{warningbox}{La memoria de parámetros no es una base de datos universal}
LoRA es adecuado para sedimentar preferencias de conducta y patrones generalizables, pero no para conservar con precisión hechos que deban poder eliminarse, auditarse o citarse uno por uno. Los registros factuales deben seguir residiendo en un memory store externo; la actualización de parámetros y la memoria recuperable deben coordinarse, no sustituirse entre sí.
\end{warningbox}

\subsection{De los adapters fuera de línea al Context Learning en línea}

Un sistema de personalización debe procesar un flujo continuo de pocas muestras a la vez, lo que exige un entrenamiento estable, de bajo costo y capaz de evitar el olvido catastrófico. El informe estudia, mediante la inicialización de LoRA, el barrido de rank/scale y protocolos de evaluación, cómo escribir la experiencia con rapidez sin dejar de conservar la capacidad de base.

\sourcefigure{0.93}{lecture06/slides-images/slide-016-00795s.jpg}{Hiperparámetros de LoRA y su efecto en la escritura de experiencia}{00:13:15--00:16:00}

\sourcefigure{0.93}{lecture06/slides-images/slide-018-01050s.jpg}{Efecto de distintas inicializaciones de LoRA en la estabilidad y la convergencia}{00:17:30--00:19:00}

\sourcefigure{0.93}{lecture06/slides-images/slide-020-01245s.jpg}{El ciclo cerrado del sistema de Context Learning: experiencia, entrenamiento y servicio}{00:20:45--00:22:45}

\begin{knowledgebox}{Escribir experiencia requiere tres puertas}
La primera es el juicio de valor: si esta interacción contiene información reutilizable en el futuro; la segunda es la atribución: si se debe actualizar la preferencia del usuario, la habilidad de la tarea o el conocimiento del entorno; la tercera es la verificación: si, tras la actualización, mejora la tarea objetivo sin dañar las capacidades previas.
\end{knowledgebox}

\subsection{MiNT: entrenamiento y servicio de modelos personalizados a escala de millones}

Cuando cada usuario puede tener su propio adapter, el problema deja de ser algorítmico y se vuelve un problema de sistemas: cómo programar el entrenamiento, mantener en caché los modelos activos, combinar lotes, y gestionar versiones y reversiones. Infraestructuras del tipo MiNT abstraen el entrenamiento y el serving de adapters como un servicio, de modo que una gran cantidad de modelos pequeños comparten los pesos base y la GPU.

\sourcefigure{0.93}{lecture06/slides-images/slide-023-01530s.jpg}{El objetivo de MiNT: entrenamiento y servicio para adapters de LLM a escala de millones}{00:25:30--00:26:30}

\sourcefigure{0.93}{lecture06/slides-images/slide-025-01605s.jpg}{Arquitectura de entrenamiento y serving de MiNT}{00:26:45--00:27:15}

\sourcefigure{0.93}{lecture06/slides-images/slide-026-01635s.jpg}{Quickstart e interfaz de servicio}{00:27:15--00:29:30}

\subsection{Del sistema de investigación al Personal Agent}

El objetivo final es combinar la memoria, los modelos personalizados y las herramientas del entorno en un Personal Agent. La inteligencia personal no se manifiesta solo en «recordar el nombre», sino en la capacidad de entender de manera sostenida el objetivo, la forma de colaborar, la tolerancia al riesgo y los proyectos en curso, y de reutilizar esa experiencia entre distintas tareas.

\sourcefigure{0.93}{lecture06/slides-images/slide-031-01830s.jpg}{El Personal Agent conecta memory, environment y el modelo personalizado}{00:30:30--00:32:15}

\begin{warningbox}{La personalización debe incluir el olvido y la autorización}
El sistema debe permitir al usuario ver, corregir, eliminar y congelar experiencias; la información sensible no debe perder capacidad de gobernanza por el solo hecho de haberse escrito en un adapter. Un producto de aprendizaje de largo plazo debe tratar el consent, la retention policy y la auditoría de versiones del modelo como funciones de primer nivel.
\end{warningbox}

\subsection{Resumen de la sección}

Context Learning convierte «qué contexto darle al modelo» en una estrategia que se puede aprender, y usa actualizaciones eficientes en parámetros para portar experiencia generalizable. La inteligencia experiencial verdaderamente utilizable depende además del entrenamiento y servicio de adapters a gran escala, de una memoria externa auditable, y de mecanismos claros de escritura, verificación y olvido.
